\documentclass[11pt,a4paper]{article}
\usepackage[utf8]{inputenc}
\usepackage[T1]{fontenc}
\usepackage{geometry}
\geometry{margin=1in}
\usepackage{hyperref}
\usepackage{enumitem}
\usepackage{booktabs}
\usepackage{longtable}
\usepackage{xcolor}
\usepackage{titlesec}

\titleformat{\section}{\Large\bfseries\color{blue!70!black}}{}{0em}{}
\titleformat{\subsection}{\large\bfseries\color{blue!50!black}}{}{0em}{}
\titleformat{\subsubsection}{\normalsize\bfseries\color{blue!30!black}}{}{0em}{}

\title{\textbf{Replication Plan}\\[0.5em]
\large Mathematical Foundations of the GraphBLAS\\[0.3em]
\normalsize OSTI ID: 1379592}
\author{Auto-generated Replication Blueprint}
\date{\today}

\begin{document}
\maketitle
\tableofcontents
\newpage

%---------------------------------------------------------------------
\section{Paper Overview}
%---------------------------------------------------------------------
This paper by Kepner et al.\ presents the mathematical foundations underlying the GraphBLAS standard---a building-block approach to graph algorithms expressed via linear-algebraic operations on generalized semirings. The work defines generalized matrix and vector operations (e.g., matrix multiply on arbitrary semirings, element-wise operations, masking, extraction, assignment) and demonstrates how classical graph algorithms (BFS, shortest paths, centrality, triangle counting, etc.) map onto these primitives.

The paper is primarily \emph{definitional/algorithmic}: it specifies algebraic structures (monoids, semirings), generalized matrix operations, and gives small worked examples. Replication therefore focuses on implementing the operations, verifying them on the paper's examples, and demonstrating the graph algorithms.

%---------------------------------------------------------------------
\section{Detailed Workflow}
%---------------------------------------------------------------------

\subsection{Phase 1: Algebraic Foundations}
\begin{enumerate}[label=\textbf{1.\arabic*}]
  \item \textbf{Define algebraic structures.}
    \begin{itemize}
      \item Implement a \texttt{Monoid} class with an associative binary operator and identity element.
      \item Implement a \texttt{Semiring} class combining an additive monoid and a multiplicative monoid satisfying distributivity; include an annihilator (zero) element.
      \item Provide built-in instances: arithmetic semiring $(\mathbb{R}, +, \times, 0, 1)$, tropical (min-plus) semiring $(\mathbb{R}\cup\{+\infty\}, \min, +, +\infty, 0)$, Boolean semiring $(\{\text{true},\text{false}\}, \lor, \land, \text{false}, \text{true})$, and others referenced in the paper.
    \end{itemize}
  \item \textbf{Define sparse matrix/vector representations.}
    \begin{itemize}
      \item Represent matrices as dictionaries of (row, col) $\to$ value or via SciPy sparse formats.
      \item Support an explicit ``no-value'' (structural zero) distinct from the additive identity when needed.
    \end{itemize}
\end{enumerate}

\subsection{Phase 2: Core GraphBLAS Operations}
\begin{enumerate}[label=\textbf{2.\arabic*}]
  \item \textbf{Matrix--matrix multiply (\texttt{mxm}).} Generalized $\mathbf{C} = \mathbf{A} \oplus.\otimes\; \mathbf{B}$ over an arbitrary semiring.
  \item \textbf{Matrix--vector multiply (\texttt{mxv}) and vector--matrix multiply (\texttt{vxm}).}
  \item \textbf{Element-wise operations (\texttt{eWiseAdd}, \texttt{eWiseMult}).} Union and intersection semantics on the structure of two matrices/vectors.
  \item \textbf{Extract and Assign sub-matrices/sub-vectors} using index arrays.
  \item \textbf{Apply} a unary function element-wise to a matrix or vector.
  \item \textbf{Reduce} a matrix to a vector (row-wise or column-wise) or a vector to a scalar using a monoid.
  \item \textbf{Transpose} operation.
  \item \textbf{Masking and accumulation}: structural and value-based masks; accumulator semantics.
  \item \textbf{Kronecker product} over a semiring.
\end{enumerate}

\subsection{Phase 3: Verification with Paper Examples}
\begin{enumerate}[label=\textbf{3.\arabic*}]
  \item Reproduce every small numeric example in the paper (e.g., the $4\times 4$ adjacency-matrix BFS example, the shortest-path tropical semiring example, the triangle-counting example).
  \item Compare outputs element-by-element against the paper's stated results.
\end{enumerate}

\subsection{Phase 4: Graph Algorithm Demonstrations}
\begin{enumerate}[label=\textbf{4.\arabic*}]
  \item \textbf{Breadth-First Search (BFS)} via repeated matrix--vector multiplication on the Boolean semiring.
  \item \textbf{Single-Source Shortest Paths} via repeated multiplication on the tropical semiring.
  \item \textbf{Triangle counting} via masked matrix multiply and reduce.
  \item \textbf{Betweenness centrality} via the BC algorithm expressed in GraphBLAS primitives.
  \item \textbf{PageRank / k-Truss / connected components} if time permits, as extended demonstrations.
\end{enumerate}

\subsection{Phase 5: Performance and Scalability (Optional)}
\begin{enumerate}[label=\textbf{5.\arabic*}]
  \item Benchmark the Python/SciPy prototype against the reference SuiteSparse:GraphBLAS C implementation on standard graph datasets (e.g., SNAP graphs, Graph500 RMAT).
  \item Report timings and memory usage.
\end{enumerate}

%---------------------------------------------------------------------
\section{Datasets}
%---------------------------------------------------------------------
The paper is primarily definitional and uses small, synthetic, hand-constructed matrices. No external datasets are required for core replication.

\begin{longtable}{p{4.5cm} p{3.5cm} p{6cm}}
\toprule
\textbf{Dataset / Source} & \textbf{Type} & \textbf{Location / Notes} \\
\midrule
\endfirsthead
\toprule
\textbf{Dataset / Source} & \textbf{Type} & \textbf{Location / Notes} \\
\midrule
\endhead
Paper's inline examples (Tables/Figures) & Synthetic small matrices & Extracted directly from the paper text \\
\midrule
SNAP graph datasets (optional benchmarking) & Real-world graphs & \url{https://snap.stanford.edu/data/} \\
\midrule
SuiteSparse Matrix Collection (optional) & Sparse matrices & \url{https://sparse.tamu.edu/} \\
\midrule
Graph500 RMAT generator (optional) & Synthetic power-law graphs & \url{https://graph500.org/} \\
\bottomrule
\end{longtable}

%---------------------------------------------------------------------
\section{Models, Tools, and Software}
%---------------------------------------------------------------------

\subsection{Programming Languages}
\begin{itemize}
  \item \textbf{Python 3.8+} --- primary implementation language for the prototype.
  \item \textbf{C/C++} --- only needed if benchmarking against SuiteSparse:GraphBLAS.
\end{itemize}

\subsection{Libraries and Frameworks}
\begin{longtable}{p{4.5cm} p{2.5cm} p{6.5cm}}
\toprule
\textbf{Tool / Library} & \textbf{Version} & \textbf{Purpose / URL} \\
\midrule
\endfirsthead
\toprule
\textbf{Tool / Library} & \textbf{Version} & \textbf{Purpose / URL} \\
\midrule
\endhead
NumPy & $\geq$ 1.21 & Dense array operations; semiring arithmetic \\
\midrule
SciPy (scipy.sparse) & $\geq$ 1.7 & Sparse matrix storage and basic linear algebra \\
\midrule
python-graphblas & latest & Python wrapper for SuiteSparse:GraphBLAS; provides high-level GraphBLAS API \newline \url{https://github.com/python-graphblas/python-graphblas} \\
\midrule
SuiteSparse:GraphBLAS & $\geq$ 7.0 & Reference C implementation of the GraphBLAS spec \newline \url{https://github.com/DrTimothyAldenDavis/GraphBLAS} \\
\midrule
NetworkX & $\geq$ 2.6 & Graph construction, visualization, and validation of algorithm outputs \\
\midrule
Matplotlib & $\geq$ 3.4 & Plotting and visualization \\
\midrule
pytest & $\geq$ 6.0 & Unit testing framework for verifying operations \\
\bottomrule
\end{longtable}

\subsection{Hardware Requirements}
\begin{itemize}
  \item For the paper's small examples: any modern laptop suffices.
  \item For optional large-scale benchmarking: multi-core workstation with $\geq$ 32\,GB RAM recommended.
\end{itemize}

\subsection{Reference Documents}
\begin{itemize}
  \item The GraphBLAS C API Specification: \url{https://graphblas.org/docs/GraphBLAS_API_C_v2.0.0.pdf}
  \item GraphBLAS Forum: \url{https://graphblas.org/}
\end{itemize}

\end{document}
